\section{Symmetric-group degrees: the quantitative toolkit}
\label{l:degrees}\label{l:degrees-section}

The lifting estimates involve sums of negative powers of dimensions.
The short proof in Lemma~\ref{l:degree-basic} already supplies every
positive inverse power. Here we develop the quantitative estimates used
by the later transfers: tableau and hook constructions, exponential
growth in the deficit from a longest line, and sharper bounds near a
row or column.

Theorem~\ref{l:degree-all-powers} combines summability with the uniform
deficit bound. Lemmas~\ref{l:degree-inverse-first}
and~\ref{l:degree-type-absorption} then control near-row dimensions and
the number of restriction types. Appendix~\ref{l:degree-supplement}
retains the independent fixed-power proofs, including routes that use
only coarse partition counts or avoid the hook formula.

For a partition $\lambda\vdash n$, write $D_\lambda=f^\lambda$ for
the dimension of its complex irreducible representation and
$\lambda'$ for its transpose.  We use the classical standard-tableau
dimension theorem and hook-length formula \cite{sagan,etingof}:
\[
 D_\lambda=\#\{\text{standard tableaux of shape }\lambda\}
 =\frac{n!}{\prod_{x\in\lambda}h(x)}.
\]
We set $D_{\varnothing}=1$ and $\binom a{-1}=0$.
A standard tableau increases along rows and columns; $h(x)$ is one
plus the number of boxes strictly to the right of or below $x$.
Transposition preserves dimension.  Put
\[
 L(\lambda)=\max(\lambda_1,\lambda'_1),\qquad
 k(\lambda)=n-L(\lambda),\qquad
 Z_u(n)=\sum_{\lambda\vdash n,\ \lambda\notin\{(n),(1^n)\}}
                     D_\lambda^{-u}.
\]
The excluded diagrams are a set, so coincide when $n=1$.  They are
exactly the one-dimensional types: transpositions are conjugate and
generate $S_n$, so a one-dimensional character takes the same value
$1$ or $-1$ on all transpositions.

\subsection{Counting constructions}

Let $p(j)$ be the number of partitions of $j$, with $p(0)=1$.
We will use each of the elementary bounds
\begin{equation}\label{l:partition-counts}
 p(j)\le2^{j-1}\le2^j,\qquad
 p(j)\le(j+1)^{2\lceil\sqrt j\rceil},\qquad
 p(j)\le e^{3\sqrt j}\quad(j\ge1).
\end{equation}
The first follows by counting compositions.  For the second, record
the multiplicities of parts at most $\lceil\sqrt j\rceil$ and a
zero-padded list of the at most $\sqrt j$ larger parts.  Each entry
has at most $j+1$ possibilities.  For the third, evaluate the partition
generating product at $x=e^{-1/\sqrt j}$.  Its logarithm is
\[
 \sum_{r\ge1}\frac1{r(e^{r/\sqrt j}-1)}
 \le\sqrt j\sum_{r\ge1}r^{-2}\le2\sqrt j,
\]
and extracting the coefficient of $x^j$ costs $e^{\sqrt j}$.
At a fixed positive value of $k(\lambda)$ there are at most $2p(k)$
shapes, since deleting a longest row or column leaves a partition of
$k$.  Also $L(\lambda)\ge\sqrt n$, because the diagram is contained
in a square of side $L(\lambda)$.

\begin{lemma}[Tableau constructions]\label{l:tableau-constructions}
Suppose a diagram has first row of length $a$, lower diagram $\gamma$
of size $k$, and second row of length $c$ (put $c=0$ if $k=0$).
\begin{enumerate}
\item Every tableau on a Young subdiagram extends to the whole diagram.
For $0\le b\le\min(a,k)$,
\begin{equation}\label{l:tableau-interleave}
 D_\lambda\ge\binom ab,
 \qquad D_\lambda\ge\binom{a+k-c}{k}.
\end{equation}
\item If $k\le a$, ballot interleavings give
\begin{equation}\label{l:tableau-ballot}
 D_\lambda\ge\binom{a+k}{k}-\binom{a+k}{k-1}
 =\binom{a+k}{k}\frac{a-k+1}{a+1}
 \ge\frac1{k+1}\binom{2k}{k}.
\end{equation}
In particular a two-row rectangle of width $b$ has
$\binom{2b}{b}/(b+1)$ tableaux.
\item If the nonempty diagonals of constant row-plus-column index
have sizes $t_1,\ldots,t_q$, then
\begin{equation}\label{l:tableau-diagonal}
 D_\lambda\ge\prod_{j=1}^q t_j!
 \ge2^{n-q},\qquad
 D_\lambda\ge\left(\frac n{qe}\right)^n.
\end{equation}
Thus width and height at most $r$ give $q\le2r-1$ and
$D_\lambda\ge2^{n-2r+1}$.
\end{enumerate}
\end{lemma}
\begin{proof}
A Young subdiagram is closed under moving up and left.  Add its
remaining boxes in row-major order with successive larger labels;
every predecessor is already present.  To prove the first bound,
retain the whole first row and a Young subdiagram of $b$ lower boxes.
Such a subdiagram is obtained by removing lower corners.  Fill the
first $b$ top boxes with $1,\ldots,b$; choose the $b$ lower labels
from the remaining $a$ labels and fill them in the relative order of
one fixed standard tableau.  All lower boxes lie in the first $b$
columns, so filling the rest of the top row increasingly is valid.
For the second bound, reserve only the first $c$ top labels and
choose all $k$ lower labels from the remaining $a+k-c$ labels.

For the ballot bound, fix a standard order of the lower boxes and
interleave their additions with those of the top row.  Require each
prefix to have at least as many top-row additions as lower additions.
The $j$th lower box lies in a column at most $j$, so its top predecessor
has been inserted.  Reflection at the first violating prefix counts
the valid words by the displayed binomial difference.  A valid
balanced word with $k$ additions of each kind, followed by $a-k$
top additions, gives the last bound.  For an actual two-row rectangle
the word condition is also necessary, proving the exact count.

Finally, put successive batches of labels on increasing diagonals,
in arbitrary order within each diagonal.  Every row or column
predecessor lies on an earlier diagonal.  This gives the product of
factorials.  The first estimate uses $t!\ge2^{t-1}$; the second uses
$t!\ge(t/e)^t$ and convexity of $t\log t$ to obtain
$\sum_jt_j\log t_j\ge n\log(n/q)$.
\end{proof}

\begin{lemma}[Separating the first-row hooks]\label{l:degree-hook-ratio}
With the notation of Lemma~\ref{l:tableau-constructions}, let $b_j$
be the number of lower boxes in column $j$.  Then
\begin{align}
 D_\lambda
 &=\binom{a+k}{k}D_\gamma
       \prod_{j=1}^a\left(1+\frac{b_j}{a-j+1}\right)^{-1}
       \label{l:hook-exact}\\
 &\ge\binom{a+k}{k}D_\gamma
       \exp\left(-\frac{k(1+\log a)}a\right).
       \label{l:hook-rearrangement}
\end{align}
If $k\le a$, a second lower bound is
\begin{equation}\label{l:hook-short-tail}
 D_\lambda\ge\binom{a+k}{k}D_\gamma
                \exp\left(-\frac{k}{a-k+1}\right).
\end{equation}
\end{lemma}
\begin{proof}
The lower hooks are unchanged, and the top hooks are $a-j+1+b_j$,
which proves the identity.  The sequence $b_j$ decreases while
$(a-j+1)^{-1}$ increases.  The opposite-order rearrangement inequality
therefore gives
\[
 \sum_{j=1}^a\frac{b_j}{a-j+1}
 \le\frac{k}{a}\sum_{j=1}^a\frac1j
 \le\frac{k(1+\log a)}a.
\]
For completeness, the rearrangement inequality follows by expanding
$\sum_{i,j}(b_i-b_j)(c_i-c_j)\le0$ for oppositely ordered
sequences $b_i,c_i$.  Apply $\log(1+x)\le x$ to the product.
For the last estimate, $b_j=0$ for $j>k$, so every relevant
denominator is at least $a-k+1$, and $\sum_jb_j=k$.
\end{proof}

\subsection{All positive powers and the first-row deficit}

\begin{theorem}[Every positive inverse power]\label{l:degree-all-powers}
For every fixed real $u>0$,
\begin{equation}\label{l:degree-all-powers-equation}
 C_u:=\sup_{n\ge1}\sum_{\lambda\vdash n}D_\lambda^{-u}<\infty,
 \qquad Z_u(n)\longrightarrow0.
\end{equation}
Moreover there are absolute constants $c_1,c_2>0$ such that, for all
sufficiently large $n$ and every $\lambda\vdash n$,
\begin{equation}\label{l:degree-tail-exponential}
 D_\lambda\ge c_1e^{c_2k(\lambda)}.
\end{equation}
For each fixed positive $k$, these dimensions tend uniformly to infinity
among shapes with $k(\lambda)=k$.
\end{theorem}
\begin{proof}
Put $L=L(\lambda)$ and $k=n-L$.  If $L\le n/16$, every hook is at
most $2L$, so
$D_\lambda\ge(n/(2eL))^n\ge2^n$.
Otherwise orient the diagram with first row $L$, and for $k\ge1$
retain $b=\min(k,\lfloor L/3\rfloor)$ lower boxes.  For sufficiently
large $n$, $b\ge1$, and Lemma~\ref{l:tableau-constructions} gives
\[
 D_\lambda\ge\binom Lb\ge(L/b)^b\ge3^b,
 \qquad b\ge k/48-1.
\]
Indeed $L>n/16$ implies $\lfloor L/3\rfloor>n/48-1\ge k/48-1$.
Thus one may take $c_1=1/3$ and $c_2=(\log3)/48$ after enlarging
the size threshold.  For a fixed positive $k$, eventually $b=k$,
and $D_\lambda\ge\binom{n-k}{k}\to\infty$ uniformly.

There are at most $2p(k)$ shapes at deficit $k$.  By
\eqref{l:partition-counts} and \eqref{l:degree-tail-exponential}, their
total inverse $u$th powers are bounded by a constant depending on $u$
times $e^{3\sqrt k-uc_2k}$.  This is summable in $k$, and each fixed
positive-$k$ contribution tends to zero.  Dominated convergence proves
$Z_u(n)\to0$.  The full sum tends to two; its finitely many initial
values are finite, including the value one at $n=1$.
\end{proof}

There are two other ways to obtain the full positive-power range.
The next proof uses ballot words; the short proof from
Lemma~\ref{l:degree-basic} uses a central-binomial subdiagram.

\begin{proof}[A ballot proof of the summability assertions]
For $k\ge1$ and $L=n-k\ge n/2$, \eqref{l:tableau-ballot} gives
\[
 D_\lambda\ge\binom nk\frac{n-2k+1}{n-k+1}
       \ge\frac{2^k}{k+1}.
\]
For the second inequality, the ballot fraction is at least $1/(k+1)$
and $\binom nk\ge(n/k)^k\ge2^k$.
The majorant $2p(k)(k+1)^u2^{-uk}$ is summable, and the first
binomial expression diverges for each fixed $k\ge1$.
If $L<n/2$ but $L\ge n/10$, retain a row of length $L$ and $L$
lower boxes.  The ballot count $\binom{2L}L/(L+1)$ is exponential
in $n$.  If $L<n/10$, the hook formula gives $(5/e)^n$.
The subexponential bound for $p(n)$ handles both remaining ranges.
Adding the one-dimensional terms proves the uniform bound.
\end{proof}

The third all-positive-power route is the one already proved in
Lemma~\ref{l:degree-basic}.  Its three ranges are $1\le k\le n/3$,
$k>n/3$ with $L\ge n/10$, and $L<n/10$.  The respective bounds are
$\binom{n-k}k\ge2^k$, $2^L/(L+1)$, and $(5/e)^n$.
Thus it uses the same initial-row construction with a central-binomial
subdiagram, rather than the deficit-exponential or ballot estimates.


\begin{lemma}[Near-row dimensions and inverse-first summability]
\label{l:degree-inverse-first}
For $1\le k=k(\lambda)\le n/4$, with $\gamma$ the lower diagram
after orienting a longest line as the first row,
\begin{equation}\label{l:degree-first-row-bound}
 D_\lambda\ge e^{-1/2}\binom nkD_\gamma
                  \ge\tfrac12\binom nk.
\end{equation}
For $k>n/4$ and all sufficiently large $n$,
$D_\lambda\ge e^{n/100}$.  These estimates independently give
$Z_1(n)\to0$ and $C_1<\infty$.
\end{lemma}
\begin{proof}
Equation~\eqref{l:hook-short-tail} has exponent
$k/(n-2k+1)\le1/2$, proving \eqref{l:degree-first-row-bound}.
If $k>n/4$ and $L\ge n/8$, keep the first row and
$b=\lfloor n/32\rfloor$ lower boxes.  For large $n$, $1\le b\le L/4$.
The same hook bound gives dimension at least
$\frac12\binom{L+b}b\ge4^b/2$.
If $L<n/8$, every hook is shorter than $n/4$, giving $(4/e)^n$.
Both exceed $e^{n/100}$ for sufficiently large $n$.
The near-row inverse sum is bounded by
$4\sum_{1\le k\le n/4}p(k)/\binom nk$.
Each term vanishes at fixed $k$ and is bounded by $4p(k)4^{-k}$;
the remaining contribution is at most $p(n)e^{-n/100}=o(1)$.
This proves both assertions about inverse-first sums.
\end{proof}

\begin{lemma}[Absorbing the number of tail types]\label{l:degree-type-absorption}
Let $M(k)=\sum_{j=0}^kp(j)$.  For every fixed $B\ge1$ and all
sufficiently large $n$, uniformly over $\lambda\vdash n$ with
$k(\lambda)\ge1$,
\begin{equation}\label{l:degree-type-absorption-equation}
 4B M(k(\lambda))\le D_\lambda^{1/4}.
\end{equation}
\end{lemma}
\begin{proof}
The partition bounds give $\log M(k)=O(\sqrt k\log(k+1))=o(k)$.
For $k>n/4$, Lemma~\ref{l:degree-inverse-first} gives
$\log D_\lambda\ge n/100$, whereas $\log M(k)=o(n)$ uniformly
for $k\le n$.  For $k\le n/4$, use
$D_\lambda\ge\frac12\binom nk\ge\frac12 4^k$ once $k$ exceeds
a fixed sufficiently large cutoff.  The finitely many positive $k$
below that cutoff are handled by the uniform divergence of
$\binom nk$.  One size threshold works for all shapes.
\end{proof}
